\documentclass[11pt,a4paper]{ctexart}
\usepackage[margin=2.2cm]{geometry}
\usepackage{xcolor}
\usepackage{enumitem}
\usepackage{amssymb}
\usepackage{booktabs}
\usepackage[normalem]{ulem}
\usepackage{array}
\usepackage{longtable}
\usepackage{titlesec}
\usepackage{hyperref}
\definecolor{navy}{HTML}{1A365D}
\hypersetup{colorlinks=true,linkcolor=navy,urlcolor=navy}
\setlist{nosep}
\titleformat{\section}{\Large\bfseries\color{navy}}{}{0pt}{}
\titleformat{\subsection}{\large\bfseries\color{navy}}{}{0pt}{}
\pagestyle{plain}
\renewcommand{\arraystretch}{1.25}

\title{\textbf{八年级下册语法总结}\\[4pt]
\large 人教版（新课纲版）八年级下册 · 按教材书末 Grammar 附录整理}
\author{初中英语资料站}
\date{2026-11-12 · 附配套专项练习 192 题}
\begin{document}
\maketitle
\section*{使用说明}
\begin{itemize}[leftmargin=1.4em]
\item 本讲义按八年级下册教材书末 Grammar 附录整理，表格与例句取自教材原文。
\item 网页版：\url{https://1901151458-prog.github.io/english-materials/g8b-grammar/}
\end{itemize}

\section{反身代词（Reflexive Pronouns）}
\noindent\textcolor{gray}{表达「……自己；本身」，由 -self / -selves 构成}\par\smallskip
反身代词是英语中表达「……自己；本身」含义的一类代词，是在第一人称、第二人称形容词性物主代词或第三人称人称代词宾格形式的词尾加 \textbf{-self} 或 \textbf{-selves} 组成。

{\footnotesize
\begin{longtable}{p{3.35cm}p{3.35cm}p{3.35cm}p{3.35cm}}
\hline
数 & 第一人称 & 第二人称 & 第三人称 \\
\hline
\endfirsthead
\hline
数 & 第一人称 & 第二人称 & 第三人称 \\
\hline
\endhead
\hline
\endfoot
单数 & myself & yourself & himself / herself / itself \\
复数 & ourselves & yourselves & themselves \\
\end{longtable}
}

{\footnotesize
\begin{longtable}{p{6.70cm}p{6.70cm}}
\hline
句法作用 & 例句 \\
\hline
\endfirsthead
\hline
句法作用 & 例句 \\
\hline
\endhead
\hline
\endfoot
动词或介词的宾语 & Be careful when you’re using the knife! Don’t cut \textbf{yourself}.\\ In teaching others, we teach \textbf{ourselves}.\\ I’ll be away for a week. You and your sister need to take care of \textbf{yourselves}.\\ Don’t be too hard on \textbf{yourself}! \\
主语或宾语的同位语 & I cooked the dinner \textbf{myself}.\\ You had better ask the dentist \textbf{herself}. \\
表语 & Just be \textbf{yourself}.\\ Mary hasn’t been quite \textbf{herself} recently. \\
\end{longtable}
}

\begin{center}\fbox{\parbox{0.92\textwidth}{\small \textbf{教材提示：}反身代词用作主语或宾语的同位语时，一般用于强调，常译为「本人；亲自」。常见后接反身代词作宾语的动词或动词短语有 \textbf{cut、hurt、protect、teach、enjoy、help、take care of、look after} 等。}}\end{center}

\begin{itemize}[leftmargin=1.4em]
\item Help \underline{yourself} to some fruit.\\ 随便吃点水果吧。（help oneself 固定搭配。）
\item The children enjoyed \underline{themselves} at the party.\\ 孩子们在聚会上玩得很开心。（enjoy oneself ＝ have a good time。）
\end{itemize}

\section{情态动词（Modal Verbs）}
\noindent\textcolor{gray}{表示提议或建议：should / could}\par\smallskip
情态动词可以用来表示提议或建议。教材重点讲了两个：\textbf{should} 和 \textbf{could}。

\subsection*{1. should}

情态动词 should 作「应该；应当；可以」讲，用以表达职责和义务、提出劝告，而且表述的是自己的主观看法，其后接\textbf{动词原形}。

\begin{itemize}[leftmargin=1.4em]
\item When you have the flu, you \underline{should wear} a mask when you’re outside.\\ 患流感时，你出门应戴上口罩。（should + 动词原形。）
\item If you did something wrong, you \underline{should say} sorry.\\ 如果你做错了什么，你应该道歉。（提出劝告。）
\end{itemize}

should 的否定式是 \textbf{should not / shouldn’t}，意思是「不应当；不该」。

\begin{itemize}[leftmargin=1.4em]
\item You \underline{shouldn’t eat} so much next time.\\ 下次你可不要吃这么多了。
\item We \underline{shouldn’t judge} a book by its cover.\\ 我们不应该以貌取人。
\end{itemize}

\subsection*{2. could}

情态动词 could 也用于提出建议，比 can 语气更为委婉。

A: What should I take with me to the party?\\ B: You \underline{could take} some fruit or flowers.\\ A：我去参加聚会应该带点什么呢？ B：你可以带些水果或鲜花。

If you want to say sorry, you \underline{could call} him.\\ 如果你想说道歉，你可以给他打电话。

\section{过去进行时（Past Continuous Tense）}
\noindent\textcolor{gray}{was / were + 动词-ing，表示过去某时正在进行}\par\smallskip
\subsection*{1. 构成和基本用法}

过去进行时由「\textbf{was / were + 动词 -ing}」构成，表示在\textbf{过去某一时刻或某一段时间正在进行的动作}。这一特定的过去时间除有上下文暗示，一般用时间状语来表示。

\begin{itemize}[leftmargin=1.4em]
\item A: What \underline{was Kate doing} at the time of the earthquake?\\ B: She \underline{was doing} her homework.\\ 地震发生时凯特在做什么？她在做作业。
\item A: What \underline{were your family doing} when the typhoon hit?\\ B: My children and I \underline{were reading} together.\\ 台风来袭时你们一家在做什么？我和孩子们在读书。
\end{itemize}

过去进行时还可用来为\textbf{一个或一系列动作的发生提供背景}。

Everyone \underline{was enjoying} the beautiful weather that morning. Some people \underline{were swimming} in the sea, and some \underline{were relaxing} on the beach.\\ 那天早上，大家都在享受美好的天气。有人在海里游泳，有人在沙滩上放松。

过去进行时既可用在主句中，也可用在从句中，还可以同时用在主句和从句中。

\begin{itemize}[leftmargin=1.4em]
\item We \underline{were visiting} our relatives when the strong winds started.\\ 强风刮起时我们正在走亲戚。
\item I \underline{was doing} some chores while my mother \underline{was preparing} lunch.\\ 妈妈在做午饭时我在做家务。
\end{itemize}

\begin{center}\fbox{\parbox{0.92\textwidth}{\small \textbf{教材提示：}持续时间的长的动作一般用过去进行时，持续时间的短的常用一般过去时。如果主句和从句都用过去进行时，则表示主句和从句中的动作都持续了一段时间。}}\end{center}

{\footnotesize
\begin{longtable}{p{6.70cm}p{6.70cm}}
\hline
肯定式 / 否定式 &  \\
\hline
\endfirsthead
\hline
肯定式 / 否定式 &  \\
\hline
\endhead
\hline
\endfoot
I / He / She / It \textbf{was working}.\\ We / You / They \textbf{were working}. & I / He / She / It \textbf{was not working}.\\ We / You / They \textbf{were not working}. \\
\textbf{Was I working?}\\ Yes, you were. / No, you were not (weren’t). & \textbf{Were you working?}\\ Yes, I was. / No, I was not (wasn’t). \\
\textbf{Was he / she / it working?}\\ Yes, he / she / it was.\\ No, he / she / it was not (wasn’t). & \textbf{Were we / you / they working?}\\ Yes, you / we / they were.\\ No, you / we / they were not (weren’t). \\
\end{longtable}
}

\subsection*{2. 一般过去时与过去进行时用法的比较}

一般过去时表示在过去某个时间发生过的动作或存在的状态；过去进行时则表示在过去某一时刻或某一段时间正在进行的动作。

\section{现在完成时（Present Perfect Tense）}
\noindent\textcolor{gray}{have / has + 过去分词，连接过去与现在}\par\smallskip
\subsection*{1. 构成和基本用法}

现在完成时由「\textbf{助动词 have / has + 过去分词}」构成。以动词 finish 为例：

{\footnotesize
\begin{longtable}{p{6.70cm}p{6.70cm}}
\hline
肯定式 & 否定式 \\
\hline
\endfirsthead
\hline
肯定式 & 否定式 \\
\hline
\endhead
\hline
\endfoot
I / You \textbf{have finished}.\\ He / She / It \textbf{has finished}.\\ We / You / They \textbf{have finished}. & I / You \textbf{have not (haven’t) finished}.\\ He / She / It \textbf{has not (hasn’t) finished}.\\ We / You / They \textbf{have not (haven’t) finished}. \\
\textbf{Have I / you finished?}\\ \textbf{Has he / she / it finished?}\\ \textbf{Have we / you / they finished?} & Yes, you / I have. / Yes, he / she / it has. / Yes, you / we / they have.\\ No, you / I have not (haven’t). / No, he / she / it has not (hasn’t). / No, you / we / they have not (haven’t). \\
\end{longtable}
}

{\footnotesize
\begin{longtable}{p{6.70cm}p{6.70cm}}
\hline
用法 & 例句 \\
\hline
\endfirsthead
\hline
用法 & 例句 \\
\hline
\endhead
\hline
\endfoot
表示过去发生或已经完成的动作对现在造成的影响或结果 & A: Have you chosen a book for your report yet?\\ B: Yes, I’ve picked \emph{The Final Problem}.（我不必再做选择了。）\\ A: Have you ever read \emph{The Three-Body Problem}?\\ B: No, I haven’t.（我对该书不了解。） \\
表示某一动作或状态从过去某一时间开始并一直持续到现在，可以和表示延续的时间状语连用 & She’s been a volunteer since 2022.\\ Mary has lived with her uncle since her parents died.\\ I have been busy these days. I haven’t had a day off for three months.\\ How long has he worked at the museum? \\
现在完成时可以 already、never、ever、just、before、yet 等连用 & I have already borrowed the book from the library.\\ He has never seen snow before.\\ Have you ever heard of \emph{Oliver Twist}?\\ I’ve just finished \emph{Alice’s Adventures in Wonderland} by Lewis Carroll. \\
have (has) been 表示「曾经到过某地」，have (has) gone 表示「已经去某地了」 & Helen has been to Xi’an once. She wants to go there again to learn more about the long history of this city.\\ Helen has gone to Xi’an. She’ll be back in a week. \\
\end{longtable}
}

\begin{center}\fbox{\parbox{0.92\textwidth}{\small \textbf{教材注释：}非延续性动词一般不能与表示持续时间的状语连用，但在否定结构中，可用非延续性动词的现在完成时表示持续的动作或状态。\\ I haven’t bought any food since last week. The fridge is empty now.\\ Where has Tom gone? We haven’t seen him the whole morning.}}\end{center}

\subsection*{2. 现在完成时和一般过去时的区别}

现在完成时表示过去发生的某一动作对现在造成的影响或结果，强调现在的情况，所以\textbf{它不能和表示过去的时间状语连用}，如 yesterday、last night、three weeks ago、in 1990 等。而一般过去时只表示过去的动作或状态，和现在不发生联系，它可以和上述表示过去的时间状语连用。

\section{动词不定式（Infinitives）}
\noindent\textcolor{gray}{to + 动词原形，可作宾语、目的状语、宾语补足语}\par\smallskip
\subsection*{1. 动词不定式的构成}

动词不定式的基本形式是「\textbf{to + 动词原形}」。动词不定式没有人称和数的变化，在句子中不能作谓语。动词不定式仍保留动词的语义特点，可以有自己的宾语和状语。

\subsection*{2. 动词不定式的用法}

动词不定式具有名词、形容词和副词的特征，因此可以作主语、宾语、宾语补足语、定语和状语等。教材重点呈现\textbf{作宾语、目的状语和宾语补足语}三种用法。

{\footnotesize
\begin{longtable}{p{6.70cm}p{6.70cm}}
\hline
句法作用 & 例句 \\
\hline
\endfirsthead
\hline
句法作用 & 例句 \\
\hline
\endhead
\hline
\endfoot
作宾语 & Do you want \textbf{to visit} the Great Wall?\\ I like \textbf{to play} badminton in my free time.\\ On Saturday, James decided \textbf{to make} fried chicken. \\
作目的状语 & What do you do \textbf{to pass} the time?\\ Allen ran to the living room \textbf{to get} the fire extinguisher.\\ Do not stand up \textbf{to reach} for food across the table. \\
作宾语补足语 & My instructor encouraged me \textbf{to get back up} whenever I fell.\\ My parents told me \textbf{not to push} myself too hard.\\ Doctor Cooper asked him \textbf{to be} more careful in the future. \\
\end{longtable}
}

\subsection*{3. 动词不定式的否定式}

动词不定式的否定形式由「\textbf{not + 动词不定式}」构成。

\begin{itemize}[leftmargin=1.4em]
\item I would prefer \underline{not to talk} about it.\\ 我宁愿不谈这件事。
\item "Try \underline{not to make} any noise when you eat," his mother said.\\ 他妈妈说：「吃饭时尽量不要发出声音。」
\end{itemize}

\subsection*{4. 「疑问词 + 不定式」结构}

动词不定式可以和 \textbf{what、which、how、where、when} 等连用。

\begin{itemize}[leftmargin=1.4em]
\item He didn’t know \underline{what to say} to Ella.\\ 他不知道该对埃拉说什么。
\item She asked me \underline{how to use} the washing machine.\\ 她问我怎么用洗衣机。
\item The tour guide will tell you \underline{where to go} next.\\ 导游会告诉你接下来去哪里。
\end{itemize}

\subsection*{5. 不带 to 的动词不定式}

有一些动词后用作宾语补足语的不定式通常不带 to，主要有两类：\textbf{感官动词}（see、hear、watch、feel、notice 等）和\textbf{使役动词}（let、make、have 等）。

\begin{itemize}[leftmargin=1.4em]
\item The referee saw Matt \underline{push} the other player.\\ 裁判看到马特推了另一个球员。
\item He heard someone \underline{walk} down the stairway.\\ 他听到有人走下楼梯。
\item If you have any questions, please let me \underline{know}.\\ 如果你有任何问题，请让我知道。
\item Hobbies can make you \underline{feel} good about yourself.\\ 爱好能让你对自己感觉良好。
\end{itemize}

\section{连词（Conjunctions）}
\noindent\textcolor{gray}{并列连词 vs 从属连词}\par\smallskip
用来连接单词、短语或句子的虚词叫连词。连词在句中不能单独作句子成分，一般不重读。按其性质，连词可分为两类：\textbf{并列连词}和\textbf{从属连词}。

\subsection*{1. 并列连词}

并列连词是指用以连接并列的单词、短语、从句或句子的连词。常见的并列连词有 \textbf{and、but、for、or、both ... and ...、not only ... but also ...} 等。

\begin{itemize}[leftmargin=1.4em]
\item You can fool some of the people all of the time, \underline{and} all of the people some of the time, \underline{but} you cannot fool all of the people all of the time.\\ 你可以一直欺骗一部分人，也可以在一段时间欺骗所有人，但你不可能一直欺骗所有人。
\item To be, \underline{or} not to be, that is the question.\\ 生存还是毁灭，这是个问题。
\end{itemize}

The Nile \underline{and} the Yellow River helped ancient civilizations to develop.

\subsection*{2. 从属连词}

从属连词是指用以引导复合句中从句的连词。常见的从属连词有 \textbf{when、while、as、before、after、until、if、because、although、though、than、so that、so ... that、unless、as soon as} 等。

\begin{itemize}[leftmargin=1.4em]
\item He started programming \underline{when} he was 10.\\ 他 10 岁时开始编程。
\item \underline{If} you run after two hares, you will catch neither.\\ 脚踩两只船，两头落空。
\item You are braver \underline{than} you believe, stronger \underline{than} you seem, and smarter \underline{than} you think.\\ 你比你相信的更勇敢，比你看起来更坚强，比你想象的更聪明。
\item Please give me a call \underline{as soon as} you arrive in Beijing.\\ 你一到北京就给我打电话。
\end{itemize}

\begin{center}\fbox{\parbox{0.92\textwidth}{\small \textbf{教材注释：}在同一个句子中，\textbf{though（although）与 but 不可同时使用}，\textbf{because 与 so 也不可同时使用}。\\ × Although you were so busy, but you took the time to help me.（应删去 but，或删去 although）\\ × Because Zhuge Liang is so clever and has a plan for everything, so I admire him.（应删去 because，或删去 so）}}\end{center}

\section{状语从句（Adverbial Clauses）}
\noindent\textcolor{gray}{时间、地点、原因、条件、结果、目的、让步七类}\par\smallskip
在复合句中，用作状语的从句叫状语从句。根据所表达的意思，状语从句可分为\textbf{时间、地点、原因、条件、结果、目的和让步}等类型。

{\footnotesize
\begin{longtable}{p{4.47cm}p{4.47cm}p{4.47cm}}
\hline
类型 & 引导连词 & 例句 \\
\hline
\endfirsthead
\hline
类型 & 引导连词 & 例句 \\
\hline
\endhead
\hline
\endfoot
时间 & before、after、when、while、as、as soon as、since、not ... until、whenever 等 & As the front door opened, James heard the sound of Allen singing.\\ You won’t know if you can do it until you try.\\ I like to write about my feelings in my diary whenever I feel bad or sad. \\
地点 & where、wherever 等 & Keep the dictionary where you can see it.\\ I bring a book wherever I go. \\
原因 & because、as、since 等 & Alice can’t reach the key on the table because she’s too small.\\ Chinese people love the ginkgo tree, as it is a symbol of hope and long life.\\ Since he has returned, you could now ask for his help. \\
条件 & if、unless 等 & If you don’t want to forget your resolutions, write them down.\\ Nothing is impossible unless you think it is. \\
结果 & so ... that & I was so hungry that I ate a big bowl of noodles and a big plate of dumplings. \\
目的 & so、so that 等 & We’ll sit near the front so we can hear the speaker better.\\ You can write to your parents so that they’ll know how you feel. \\
让步 & although、even though 等 & Although my guitar is old, it means a lot to me.\\ Tom felt proud of their team, even though they lost. \\
\end{longtable}
}

\begin{center}\fbox{\parbox{0.92\textwidth}{\small \textbf{教材注释：}在\textbf{时间状语从句和条件状语从句}中，通常用\textbf{一般现在时}表示一般将来时。\\ Please say hello for me when you see her next week.\\ If it doesn’t rain tomorrow, we’ll go for a picnic.}}\end{center}

\section{配套专项练习}
\noindent\textcolor{gray}{四个专项各 48 题、共 192 题，支持在线作答与即时解析}\par\smallskip
\paragraph{时态专项 · 48 题}

\paragraph{动词不定式专项 · 48 题}

\paragraph{连词与状语从句专项 · 48 题}

\paragraph{反身代词与情态动词专项 · 48 题}

\begin{center}\fbox{\parbox{0.92\textwidth}{\small \textbf{建议顺序：}先按本页 01–07 讲通规则 → 再进入对应专项练习自测 → 错题回到小节复习。四个专项合计 \textbf{192 题}。\\ \textbf{考前速查：}一页纸速查表 →（A4 单页，可打印）}}\end{center}

\end{document}
